\textbf{1. Kiến thức cần nhớ}
\begin{itemize}
\item \textbf{Đường tròn ngoại tiếp tam giác} là đường tròn đi qua ba đỉnh của tam giác.
\begin{itemize}
\item Tâm $O$ là giao điểm của ba đường \textbf{trung trực}.
\item Với tam giác vuông: Tâm là \textbf{trung điểm cạnh huyền}, bán kính $\boxed{R = \dfrac{\text{cạnh huyền}}{2}}$.
\item Với tam giác đều cạnh $a$: Tâm là trọng tâm, bán kính $\boxed{R = \dfrac{a\sqrt{3}}{3}}$.
\end{itemize}
\item \textbf{Đường tròn nội tiếp tam giác} là đường tròn tiếp xúc với ba cạnh của tam giác.
\begin{itemize}
\item Tâm $I$ là giao điểm của ba đường \textbf{phân giác trong}.
\item Bán kính $r$: $\boxed{S = p \cdot r \implies r = \dfrac{S}{p}}$ ($p$ là nửa chu vi).
\item Với tam giác đều cạnh $a$: $\boxed{r = \dfrac{a\sqrt{3}}{6}} = \dfrac{R}{2}$.
\item Với tam giác vuông tại $A$: $\boxed{r = \dfrac{b + c - a}{2}}$.
\end{itemize}
\end{itemize}

\textbf{2. Phương pháp giải các dạng}

%% DANG 9H9?28-1
\textbf{Dạng 1. Tâm, bán kính đường tròn ngoại tiếp tam giác}
\begin{enumerate}
\item Xác định dạng của tam giác (tam giác vuông, đều, cân hay thường).
\item Với tam giác vuông $\rightarrow$ dùng trung điểm cạnh huyền $R = \dfrac{BC}{2}$.
\item Với tam giác thường: dùng công thức diện tích $S = \dfrac{abc}{4R} \implies R = \dfrac{abc}{4S}$.
\end{enumerate}

%% DANG 9H9?28-2
\textbf{Dạng 2. Tâm, bán kính đường tròn nội tiếp tam giác}
\begin{enumerate}
\item Tính diện tích tam giác $S$ và nửa chu vi $p = \dfrac{a+b+c}{2}$.
\item Tính bán kính $r = \dfrac{S}{p}$.
\item Nhớ hệ thức tiếp tuyến cắt nhau từ $3$ đỉnh đến các tiếp điểm của đường tròn nội tiếp.
\end{enumerate}

%% DANG 9H9?28-3
\textbf{Dạng 3. Tam giác đều, tam giác vuông và đường tròn}
\begin{enumerate}
\item Trong tam giác đều: Tâm đường tròn ngoại tiếp và nội tiếp trùng nhau (là trọng tâm). Bán kính ngoại tiếp gấp đôi bán kính nội tiếp: $\boxed{R = 2r}$.
\item Vận dụng các công thức liên hệ nhanh để giải các bài toán thực tế hoặc tính diện tích hình vành khăn giữa hai đường tròn.
\end{enumerate}
